\chapter*{Preface}
\addcontentsline{toc}{chapter}{Preface}

Image processing is an unusually effective way to learn numerical programming.
An image is immediately recognisable, yet it can also be treated as structured
numerical data. A small change to an algorithm can often be seen at once in the
resulting image, while the same program raises deeper questions about data
representation, numerical range, file formats, colour models, filtering,
geometry and compression.

The starting point for this book is personal as well as historical. In 2001 I
took UQC146S1, \emph{Introductory Image Processing in C}, at the University of
the West of England (UWE). The module was taught by Ian Johnson, a lecturer at
UWE at the time and the module author named in the surviving specification
\cite{uwe_uqc146_spec}. The module combined a limited introduction to C with
practical image-processing work\cite{johnson_uqc146_intro}. Its course website
directed students towards progressively implementing techniques introduced in
lectures and maintaining a portfolio of image-processing functions
\cite{johnson_uqc146_web}. The coursework specification explicitly described
that portfolio as a library of functions and gave importance not only to code,
but also to programming style, documentation and testing
\cite{johnson_uqc146_assignment}.

This book preserves that implementation-led spirit while moving the work into a
modern Python environment. Algorithms are developed and understood before they
are replaced by convenient library calls. NumPy, Pillow, OpenCV and
scikit-image are valuable tools, but they become more useful once the behaviour
beneath their APIs is understood.

A reusable Python package, \pyimage{}, is developed alongside the text. The
companion Jupyter notebooks provide executable chapter material, while a pytest
suite verifies the library independently. The intention is that the code shown
in the book should form part of a working software project rather than a set of
isolated fragments.

This is a development edition. References, terminology, algorithms and
implementation details will be checked chapter by chapter against appropriate
primary and authoritative sources before the manuscript is considered ready for
publication.
